\noindent\textbf{Elección del procedimiento.} Una línea de corriente congela el tiempo: $dy/dx=at/u_0$, recta de pendiente constante. La trayectoria integra $dx/dt=u_0$ y $dy/dt=at$ con origen inicial; resulta $x=u_0t$, $y=at^2/2$. Eliminar $t=x/u_0$ da la parábola.

\begin{enumerate}

\item Para la línea de corriente se fija el instante de observación.

\[\fitmath{\mathrm{slope}=\frac{a\cdot \mathrm{tobs}}{u0}}\]

\[\fitmath{\mathrm{slope}=\frac{(0.5)\cdot (2)}{(2)}\approx 0.5\,}\]

\item La línea que pasa por el origen tiene ordenada pendiente por abscisa.

\[\fitmath{\mathrm{ystream}=\mathrm{slope}\cdot x}\]

\[\fitmath{\mathrm{ystream}=(0.5)\cdot (4)\approx 2\,\mathrm{m}}\]

\item La trayectoria usa el tiempo que tarda la partícula en llegar a esa abscisa.

\[\fitmath{\mathrm{tpart}=\frac{x}{u0}}\]

\[\fitmath{\mathrm{tpart}=\frac{(4)}{(2)}\approx 2\,\mathrm{s}}\]

\item Integrar la velocidad vertical creciente, no congelarla en el instante final.

\[\fitmath{\mathrm{ypart}=\frac{a\cdot {\mathrm{tpart}}^{2}}{2}}\]

\[\fitmath{\mathrm{ypart}=\frac{(0.5)\cdot {(2)}^{2}}{2}\approx 1\,\mathrm{m}}\]

\end{enumerate}
